\title{Direct Preference Optimization: Your Language Model is Secretly a Reward Model}

\begin{abstract}
Large-scale unsupervised language models (LMs) acquire extensive world knowledge and certain reasoning capabilities, but exerting fine-grained control over their outputs remains challenging because their training is entirely unsupervised. Methods designed to improve this controllability typically rely on collecting human preference data regarding the quality of model outputs and subsequently fine-tuning the LM to better align with those preferences, commonly employing reinforcement learning from human feedback (RLHF). Yet, RLHF is an intricate and often unstable process, requiring the initial training of a reward model to encapsulate human judgments and subsequent reinforcement learning to adjust the LM such that it maximizes the learned reward while staying proximate to the original model. In this work, we present a novel reparameterization of the reward model within RLHF, which makes it possible to derive the corresponding optimal policy analytically, thus recasting the conventional RLHF framework into one that can be solved via a straightforward classification loss. This leads to an algorithm we term \textit{Direct Preference Optimization} (DPO), which is stable, effective, and resource-efficient; it obviates the need for language model sampling during the fine-tuning process and dispenses with extensive hyperparameter searches. Our empirical studies demonstrate that DPO matches or surpasses the capability of current techniques in adapting LMs to human preferences. In particular, DPO fine-tuning outperforms PPO-based RLHF in controlling generative sentiment and delivers equivalent or superior performance in tasks such as summarization and single-turn dialogue, all while being considerably more streamlined to implement and execute.
\end{abstract}

\section{Introduction}
Large-scale unsupervised language models (LMs) trained on extensive corpora have demonstrated unexpected proficiencies~\citep{chowdhery2022palm, brown2020language, touvron2023llama,bubeck2023sparks}. Yet, the data fueling these models originates from humans with a diverse array of motives, priorities, and expertise. Not all of these characteristics are necessarily desirable for the model to reproduce; for instance, while we may want an AI coding assistant to \textit{recognize} prevalent programming errors to address them effectively, we still prefer its code generation to reflect the (possibly infrequent) instances of high-caliber coding skill within its training examples. Similarly, it is useful for a language model to be \textit{informed} of a widespread misconception held by half the population, but we emphatically do not wish for it to endorse this misconception in half its answers on the topic. Therefore, it is essential to \emph{selectively direct} the model’s output and conduct from the broad spectrum of its acquired \textit{knowledge and capacities} to ensure AI systems that are safe, effective, and controllable \citep{ouyang2022training}. Standard techniques typically align LMs with human preferences by means of reinforcement learning (RL); in this work, we demonstrate that the objective optimized in current RL approaches can be matched precisely by a simple binary cross-entropy loss, thereby streamlining the entire preference optimization process.

Broadly, the conventional workflow injects target behaviors into language models through curated collections of human preferences indicative of desirable and safe conduct. This phase of preference learning succeeds an initial unsupervised pre-training stage on a vast dataset of raw text. While a naive method would simply involve supervised fine-tuning with human-generated exemplars of preferred responses, reinforcement learning from human or artificial feedback (RLHF/RLAIF; \citep{christiano2017deep,bai2022constitutional}) has emerged as the leading paradigm. RLHF procedures train a reward model to encode human preferences, then leverage RL algorithms to adapt the language model’s outputs toward higher-reward responses while limiting deviation from the original pre-trained model. Although RLHF has led to notable gains in dialogue and coding competence, its application requires substantially more complexity compared to supervised approaches, mandating the training of multiple models and continual policy sampling during learning, thereby resulting in substantial computational overhead.

This paper presents a method for optimizing a language model to better reflect human preferences, without the need for explicit reward modeling or the use of reinforcement learning. We introduce \textit{{Direct Preference Optimization} (DPO)}, an algorithm that implicitly targets the same goal as popular RLHF techniques—namely, maximizing a reward function under a KL-divergence constraint—while remaining easy to implement and train. Conceptually, DPO works by boosting the log probability ratio in favor of preferred responses over those deemed less preferable, and employs a dynamically calculated, sample-specific importance weight. This importance weighting mitigates model collapse that we observed with more naive probability ratio objectives. As with earlier approaches, DPO assumes access to a formal preference model—such as the Bradley-Terry model \cite{bradley1952rankanalysis}—for quantifying the alignment between a reward signal and observed preference data. Distinctively, however, rather than using the preference model solely to define a loss for training a separate reward model (and subsequently a policy via that reward), DPO uses a change of variables approach so that the preference loss is expressed as a function of the policy itself. This enables direct policy optimization based on binary cross-entropy when provided with a dataset of human preference judgments over model outputs, and yields an optimal policy with respect to an implicit reward derived from those preferences.

Our primary contribution is Direct Preference Optimization (DPO), a streamlined and RL-free approach for language model training using human preference data. Experimental results demonstrate that DPO matches or outperforms established methods, including RLHF techniques based on PPO, across preference-driven tasks such as sentiment control, summarization, and dialogue, and is effective for models with up to 6B parameters.

\section{Related Work}

Progressively larger self-supervised language models are able to handle certain tasks in a zero-shot manner \citep{radford2019language} or with only a handful of examples provided as prompts \citep{gpt3,megatron,chowdhery2022palm}. Nevertheless, their effectiveness on downstream tasks and the degree to which their outputs match user expectations are substantially enhanced through fine-tuning on curated collections of instructions and human-generated completions \citep{mishra-etal-2022-cross,sanh2022multitask,chung2022scaling,thoppilan2022lamda}. This method, termed `instruction-tuning', allows large language models to better generalize beyond the explicit instructions present in the fine-tuning dataset, resulting in broader applicability and improved practical use \citep{chung2022scaling}. While instruction-tuning has demonstrated considerable effectiveness, obtaining human \textit{relative} quality ratings is often more feasible than sourcing expert-labeled demonstrations. As a consequence, subsequent research has advanced by refining LLMs with datasets consisting of human preference judgments, which in turn has led to performance gains in areas like translation \citep{kreutzer-etal-2018-reliability}, summarization \citep{stiennon2022learning,ziegler2020finetuning}, narrative generation \citep{ziegler2020finetuning}, and following instructions \citep{ouyang2022training,ramamurthy2023is}. Typically, these strategies involve first fitting a neural network-based reward model to reflect the observed preference data using models such as the Bradley-Terry preference model \citep{bradley1952rankanalysis}, and subsequently fine-tuning the language model via reinforcement learning to maximize the predicted reward—employing algorithms like REINFORCE \citep{williams1992reinforce}, proximal policy optimization (PPO; \cite{schulman2017proximal}), or related adaptations \citep{ramamurthy2023is}. Another closely connected approach utilizes instruction-following LLMs, already enhanced by human feedback, to produce additional synthetic preference datasets that focus on specific properties such as safety or harmlessness \citep{bai2022constitutional}, relying primarily on weak human supervision via text-based rubrics for the LLM's evaluations. These efforts effectively combine advancements from two areas: one focused on applying reinforcement learning for various language model objectives \citep{Ranzato2015SequenceLT,paulus2018a,wu2018learning}, and the other on learning more generally from human preferences \citep{christiano2017deep,kupcsik2018learning}. Although relative preference learning is appealing, the process of refining large language models using reinforcement learning remains challenging in practical terms. The present work proposes a theoretically grounded method for optimizing models based on relative preference data without resorting to reinforcement learning.

Research on policy learning from preferences outside of language-based tasks has been conducted in both bandit and reinforcement learning frameworks, leading to the development of various methodologies. In particular, contextual bandit learning that relies on preferences or action rankings—rather than numerical rewards—is referred to as the contextual dueling bandit (CDB) problem (\cite{yue2012karmed,dudik2015contextual}). Since absolute reward signals are not available in this setting, theoretical investigations of CDBs replace the standard notion of an optimal policy with that of a \textit{von Neumann winner}: a policy whose expected probability of winning against any alternative policy meets or exceeds 50\% \citep{dudik2015contextual}. One distinction is that, in the CDB framework, preferences are supplied in an online manner; by contrast, preference learning from human feedback often relies on a predetermined, offline dataset comprising labeled pairs of actions \citep{yan2022human}. Along similar lines, \textit{preference-based RL} (PbRL) employs binary feedback derived from an unknown ‘scoring’ function instead of explicit reward signals \citep{BusaFekete2014,ruiz2023dueling}. There exist a range of PbRL algorithms, including those that leverage off-policy preference data; however, these typically involve a two-step process of first estimating a hidden scoring (reward) function, followed by its optimization \citep{jain2013learning,BusaFekete2014,christiano2017deep,sadigh2017active,kupcsik2018learning}. In contrast, our work introduces a unified approach that learns policies in a single phase, directly maximizing alignment with preferences.

\section{Preliminaries}\label{section:prelims}

Here, we summarize the RLHF pipeline as described by \citeauthor{ziegler2020finetuning} (see also \citep{stiennon2022learning, bai2022training, ouyang2022training}), which generally consists of three main components: 1) supervised fine-tuning (SFT), 2) collecting preference data and constructing a reward model, and 3) reinforcement learning-based optimization.

\textbf{SFT}: The RLHF process initiates with supervised fine-tuning of a pre-existing language model, adapting it to the desired application (such as dialogue generation, summarization, etc.) using high-quality labeled data. This produces a refined model, $\pi^\text{SFT}$.

\textbf{Reward Modelling Phase}: In the next stage, the SFT-tuned model receives prompts $x$ and generates answer pairs $(y_1, y_2)\sim \pi^\text{SFT}(y \mid x)$. Human annotators are then tasked with selecting their preferred answer for each prompt, resulting in feedback such as $y_w\succ y_l \mid x$, where $y_w$ and $y_l$ identify the favored and non-favored responses, respectively. These preferences are hypothesized to originate from an unobserved reward function $r^*(y, x)$. To infer the structure of these preferences, models such as the Bradley-Terry (BT) model \cite{bradley1952rankanalysis} are often employed (though more general frameworks like the Plackett-Luce model \citep{plackett1975analysis, luce2012individual} are also applicable if ranked lists are available). Under the BT model, the distribution over human preferences $p^*$ is formulated as:

Given a static dataset of comparisons, denoted as $\mathcal{D}=\bigl\{x^{(i)}, y_w^{(i)}, y_l^{(i)}\bigr\}_{i=1}^N$ and presumed to be sampled from $p^*$, we introduce a parameterized reward function $r_{\phi}(x, y)$ and learn its parameters using maximum likelihood estimation. When cast as a binary classification task, the corresponding loss takes the form of the negative log-likelihood:

\begin{equation}\label{eq:reward_model}
    \mathcal{L}_R(r_{\phi}, \mathcal{D}) = -\mathbb{E}_{(x, y_w, y_l)\sim \mathcal{D}}\bigl[\log \sigma(r_{\phi}(x, y_w)- r_{\phi}(x, y_l))\bigr]
\end{equation}

where $\sigma$ denotes the logistic sigmoid function. In the case of large language models, the mapping $r_{\phi}(x, y)$ is commonly initialized from the SFT model $\pi^\text{SFT}(y \mid x)$, augmented with a final linear projection layer atop the last transformer block to yield a scalar reward estimate~\cite{ziegler2020finetuning}. To control for reward scale and reduce variance, previous works apply normalization to ensure that $\mathbb{E}_{x,y\sim \mathcal{D}}\left[r_\phi(x, y)\right] = 0$ for all $x$.

\textbf{RL Fine-Tuning Phase}: In the RL fine-tuning stage, the acquired reward function serves as the feedback mechanism for adjusting the language model. Building on earlier research~\citep{jaques2017sequence, jaques2020human}, the optimization problem is formalized as

\begin{equation}\label{eq:RL}
\max_{\pi_{\theta}}  \mathbb{E}_{x\sim \mathcal{D}, y\sim \pi_{\theta}(y \mid x)}\bigl[r_{\phi}(x, y)\bigr] - \beta\mathbb{D}_{\textrm{KL}}\bigl[\pi_{\theta}(y\mid x)\mid \mid \pi_\text{ref}(y\mid x)\bigr],
\end{equation}

where $\beta$ regulates the allowance for divergence from the reference policy $\pi_\text{ref}$, typically the initial SFT model $\pi^\text{SFT}$. The language model policy $\pi_\theta$ is usually initialized using $\pi^\text{SFT}$ as well. Imposing this KL penalty plays a crucial role, as it constrains the model to remain within the regime where the reward function is reliable, while also helping to preserve diversity in outputs and prevent collapse to a handful of high-reward responses. Because text generation operates over discrete tokens, the above objective cannot be optimized with standard gradient methods and instead relies on reinforcement learning. Common practice~\citep{ziegler2020finetuning, stiennon2022learning, bai2022training, ouyang2022training} reparameterizes the objective as ${r(x, y) = r_{\phi}(x, y) -\beta (\log \pi_{\theta}(y\mid x) - \log \pi_\text{ref}(y\mid x))}$ and applies the PPO algorithm~\cite{schulman2017proximal} to perform maximization.

\section{Direct Preference Optimization}\label{sec:DPO}

Considering the substantial obstacles posed by RL methods for large-scale tasks like language model fine-tuning, we aim to formulate a more direct approach that optimizes policies with respect to preference data, bypassing reinforcement learning altogether. Distinct from prior RLHF strategies that involve first fitting a reward function and subsequently applying RL, our methodology utilizes a particular parameterization for the reward model, which allows the optimal policy to be recovered analytically in closed form—thereby sidestepping the need for a separate RL optimization phase. As detailed below, our main contribution lies in exploiting an explicit transformation from reward functions to their corresponding optimal policies, which in turn allows us to recast the original reward-based loss as a policy-based loss. This reparameterization obviates the necessity of fitting a stand-alone reward predictor, yet still enables learning consistent with common preference modeling frameworks such as the Bradley-Terry model. Effectively, the policy model takes on the dual role of generating language and implicitly representing reward.

\textbf{DPO Objective Derivation.} We begin from the reinforcement learning objective presented in Eq.~\ref{eq:RL}, employing a general reward function $r$, as established in preceding literature. Previous work~\citep{peters2007reinforcement, peng2019advantage, korbak2022reinforcement, go2023aligning} has demonstrated that the optimal policy under the KL-regularized reward maximization in Eq.~\ref{eq:RL} has the structure:

\begin{equation}\label{eq:op_policy}
    \pi_r(y\mid x) = \frac{1}{Z(x)}\pi_\text{ref}(y\mid x)\exp\left(\frac{1}{\beta}r(x, y)\right),
\end{equation}

where the partition function $Z(x) =\sum_{y}\pi_\text{ref}(y\mid x)\exp\left(\frac{1}{\beta}r(x, y)\right)$. For a complete derivation, see Appendix \ref{app:derivation1}. Even when using $r_{\phi}$, the maximum likelihood estimator for the underlying true reward $r^*$, directly estimating $Z(x)$ remains computationally demanding~\citep{korbak2022reinforcement, go2023aligning}, limiting practical use of this form. Nevertheless, Eq.~\ref{eq:op_policy} can be manipulated to write the reward function in terms of the associated optimal policy $\pi_r$, the reference policy $\pi_\text{ref}$, and the partition function $Z(\cdot)$. Taking the logarithm of both sides of Eq.~\ref{eq:op_policy} and rearranging yields:

\begin{equation}\label{eq:main_eq}
    r(x,y) =\beta \log \frac{\pi_r(y\mid x)}{\pi_\text{ref}(y\mid x)} + \beta \log Z(x).
\end{equation}

This formulation can be directly applied to the true reward $r^*$ and its optimal policy $\pi^*$. Notably, the Bradley-Terry model depends solely on the difference in rewards between two outputs, i.e., ${p^*(y_1 \succ y_2 \mid x) = \sigma(r^*(x, y_1) - r^*(x, y_2))}$. By plugging the above reparameterization from Eq.~\ref{eq:main_eq} for $r^*(x, y)$ into the preference likelihood from Eq.~\ref{eq:bradley-terry}, the partition function term cancels out, expressing the human preference probability entirely via the optimal policy $\pi^*$ and the reference policy $\pi_\text{ref}$. Therefore, under the Bradley-Terry model, the optimal RLHF policy $\pi^*$ must satisfy:

\begin{equation}\label{eq:objective}
    p^*(y_1\succ y_2 \mid x)=\frac{1}{1 + \exp\left(\beta \log \frac{\pi^*(y_2\mid x)}{\pi_\text{ref}(y_2\mid x)} - \beta \log \frac{\pi^*(y_1\mid x)}{\pi_\text{ref}(y_1\mid x)}\right)}
\end{equation}

Further details are provided in Appendix~\ref{app:derivation2}. While this formulation assumes the Bradley-Terry model, similar derivations extend to the broader class of Plackett-Luce models~\citep{plackett1975analysis, luce2012individual}, as documented in Appendix~\ref{app:plackett_luce_models}.

Given that the probability of human preferences is now written in terms of the optimal policy, rather than in terms of the reward function, we can establish a maximum likelihood training objective for a parameterized policy $\pi_\theta$. Mirroring the reward-modeling perspective (cf. Eq.~\ref{eq:reward_model}), the resulting policy optimization problem is:

\begin{equation}\label{eq:optimum_model}
    \mathcal{L}_\text{DPO}(\pi_{\theta}; \pi_\text{ref}) = -\mathbb{E}_{(x, y_w, y_l)\sim \mathcal{D}}\left[\log \sigma \left(\beta \log \frac{\pi_{\theta}(y_w\mid x)}{\pi_\text{ref}(y_w\mid x)} - \beta \log \frac{\pi_{\theta}(y_l\mid x)}{\pi_\text{ref}(y_l\mid x)}

In this formulation, we learn an implicit reward function via an alternative parameterization, where the resulting optimal policy corresponds directly to $\pi_\theta$. Additionally, because our approach mirrors fitting a reparameterized Bradley-Terry model, it inherits several theoretical guarantees, including consistency provided appropriate conditions on the preference data distribution are met \cite{bong2022generalized}. Section~\ref{sec:theory} elaborates further on DPO’s theoretical attributes and draws comparisons with prior methods.

\textbf{Mechanistic interpretation of the DPO update.} To gain deeper insight into the behavior of DPO, we can inspect the gradient of the loss $\mathcal{L}_\text{DPO}$. With respect to the parameters $\theta$, this gradient is expressed as:

\begin{multline*}\label{eq:gradient}
    \nabla_\theta \mathcal{L}_\text{DPO}(\pi_\theta;\pi_\text{ref}) = \\ -\beta\mathbb{E}_{(x, y_w, y_l) \sim \mathcal{D}} \bigg[\underbrace{\sigma(\hat{r}_\theta(x, y_l) - \hat{r}_\theta (x, y_w))}_\text{greater weight when reward ranking is incorrect}\bigg[\underbrace{\nabla_\theta\log \pi(y_w \mid x)}_\text{encourage $y_w$ predictions} - \underbrace{\nabla_\theta\log\pi(y_l \mid x)}_\text{discourage $y_l$ predictions}\bigg]\bigg],
\end{multline*}

Here, $\hat{r}_\theta(x, y) = \beta \log \frac{\pi_\theta(y \mid x)}{\pi_\text{ref}(y \mid x)}$ defines the implicit reward associated with the candidate response as estimated by the current model $\pi_\theta$ relative to a reference $\pi_\text{ref}$ (see additional details in Section~\ref{sec:theory}). Conceptually, this loss gradient amplifies the probability of preferred outputs $y_w$, while diminishing that of less preferred outputs $y_l$. The contribution of each training triplet is modulated by how much the implicit reward model $\hat{r}_\theta$ overvalues the less preferred completion $y_l$, scaled by $\beta$, effectively reflecting the severity of the ranking error and the tightness of the KL regularization. Our empirical analyses highlight the critical role of this weighting; omitting the associated coefficient—as in a basic version of the algorithm—may lead to model collapse (see Appendix Table~\ref{tab:unlikelihood_generations}).

\textbf{DPO overview.} The DPO process generally consists of the following steps: First, completions $y_1, y_2 \sim \pi_\text{ref}(\cdot \mid x)$ are generated for each input prompt $x$, then annotated according to human preferences to form an offline preference dataset, denoted as $\mathcal{D} = \{x^{(i)}, y_w^{(i)}, y_l^{(i)}\}_{i=1}^N$. Next, the language model $\pi_\theta$ is trained by minimizing $\mathcal{L}_\text{DPO}$ with respect to a chosen reference model $\pi_\text{ref}$, dataset $\mathcal{D}$, and the specified value of $\beta$.

In practical scenarios, it is desirable to leverage existing publicly released preference datasets, instead of collecting new samples and associated human judgments. Preference data is often collected with a particular model, typically $\pi^\text{SFT}$, so we set $\pi_\text{ref} = \pi^\text{SFT}$ when such a model is accessible. When this is not an option, we estimate $\pi_\text{ref}$ by fitting it to maximize the likelihood of the preferred completions ${(x, y_w)}$, formally, ${\pi_\text{ref} = \argmax_{\pi}\mathbb{E}_{x, y_w \sim \mathcal{D}}\left[\log \pi(y_w \mid x)\right]}$. Employing this strategy helps address discrepancies between the ideal reference distribution (which is not available) and the actual $\pi_\text{ref}$ utilized during DPO training. Implementation specifics and hyperparameter choices are further described in Appendix~\ref{app:implementation}.

\section{Theoretical Analysis of DPO} This section provides a deeper examination of the DPO framework, develops theoretical justification, and connects the strengths of DPO to shortcomings of actor-critic methods used in RLHF, such as PPO~\cite{schulman2017proximal}.

\label{sec:theory}

\subsection{Your Language Model Is Secretly a Reward Model} DPO circumvents both the explicit learning of a reward model and separate reinforcement learning procedures by employing a unified maximum likelihood approach. Importantly, the optimization target in Eq. \ref{eq:main_eq} aligns with the Bradley-Terry formulation, where the reward function takes the form $r^*(x, y) = \beta \log\frac{\pi^*_\theta(y \mid x)}{\pi_\text{ref}(y \mid x)}$. Training the parameterized model $\pi_\theta$ under this objective is thus equivalent to reward model training as in Eq. \ref{eq:reward_model}, but expressed through a change of variables. This section constructs the formal foundations for such a reparameterization, establishes that it does not limit the representational power of learned reward models, and demonstrates that the optimal policy can be exactly obtained. We begin by introducing a notion of equivalence between reward functions.

\begin{definition}
Two reward functions $r(x, y)$ and $r'(x, y)$ are defined as equivalent if and only if there exists a function $f$ such that ${r(x, y)-r'(x, y) = f(x)}$.    
\end{definition}

This notion forms an equivalence relation, dividing the set of reward functions into equivalence classes. The following two lemmas capture key implications:

\begin{lemma}\label{lemma:same_prefrence} 
    Within the Plackett-Luce framework, including its special case Bradley-Terry, any two reward functions that belong to the same equivalence class will result in identical preference distributions.
\end{lemma}

\begin{lemma}\label{lemma:same_policy}
    For the constrained RL formulation, reward functions in the same equivalence class lead to the same optimal policy.
\end{lemma}

The proofs for these statements are elementary and can be found in Appendix~\ref{app:lemma1}. The first lemma describes the classic issue of under-specification inherent to the Plackett-Luce model family \cite{plackett1975analysis}. Owing to this ambiguity, it is customary to enforce further identifiability constraints to guarantee the properties of maximum likelihood estimation for Eq.~\ref{eq:reward_model} \cite{bong2022generalized}. The second lemma confirms that all representatives of an equivalence class yield the same optimal policy. Thus, for our overall aim, it suffices to identify any reward function within the optimal equivalence class. The result below, proven in Appendix~\ref{app:thm1}, formalizes this idea:

\begin{theorem}\label{thm:main}
    Provided mild technical assumptions, every equivalence class of reward functions consistent with the Plackett-Luce family (and particularly Bradley-Terry) admits a reparameterization of the form ${r(x, y) = \beta \log \frac{\pi(y\mid x)}{\pi_\text{ref}(y\mid x)}}$ for some model $\pi(y\mid x)$ and a specified reference policy $\pi_\text{ref}(y \mid x)$.
\end{theorem}

\begin{sproof}
    Let $r(x, y)$ denote a generic reward function that yields an associated optimal policy $\pi_r(y \mid x)$, as described in Eq.~\ref{eq:op_policy}. We will demonstrate that an element from the equivalence class of $r$ can always be written using the aforementioned reparameterization. Define the transformation $f$ by
\begin{equation}
    f(r; \pi_\text{ref}, \beta)(x, y) = r(x, y) - \beta\log\sum_{y}\pi_\text{ref}(y\mid x)\exp\left(\frac{1}{\beta}r(x, y)\right)
\end{equation}
Here, $f$ effectively rescales the reward function by subtracting the log-partition function (with respect to $\pi_\text{ref}$). Since this normalization only depends on $x$, $f(r; \pi_\text{ref}, \beta)(x, y)$ resides in the same equivalence class as $r(x, y)$. Moreover, substituting $r$ with the right-hand side of Eq.~\ref{eq:main_eq} (which holds generally), it follows that $f(r; \pi_\text{ref}, \beta)(x, y) = \beta \log \frac{\pi_r(y\mid x)}{\pi_\text{ref}(y\mid x)}$. Hence, the mapping $f$ recovers a representative in the desired parametric form from the equivalence class of $r$, implying no loss of generality in employing the proposed reparameterization for the reward model.
\end{sproof}

An alternative interpretation of Theorem~\ref{thm:main} is that it identifies the specific reward function chosen by the DPO reparameterization within each equivalence class, namely, the reward function that fulfills:

\begin{equation}\label{eq:lag_p}
     \sum_{y}\underbrace{\pi_\text{ref}(y\mid x)\exp\left(\frac{1}{\beta}r(x, y)\right)}_{=\pi(y\mid x)\text{, using Thm.~\ref{thm:main} reparam.}} = 1,
\end{equation}

In other words, $\pi(y\mid x)$ forms a legitimate probability distribution, as its values are non-negative and sum to one. Moreover, referencing Eq.~\ref{eq:op_policy}, it becomes evident that Eq.~\ref{eq:lag_p} represents the partition function associated with the optimal policy dictated by the reward function $r(x, y)$. The central idea underlying the DPO algorithm is to introduce explicit constraints to the otherwise underdetermined Plackett-Luce (and more specifically, Bradley-Terry) class of preference models. This preserves the set of reward functions that can be represented, while at the same time ensuring that the optimal policy described in Eq.~\ref{eq:op_policy} remains computationally manageable for any prompt $x$.

\subsection{Instability of Actor-Critic Algorithms} Our proposed framework can further elucidate sources of instability inherent in conventional actor-critic approaches for RLHF, such as PPO. Following the standard RLHF protocol, our focus is on the RL fine-tuning phase as discussed in Section~\ref{section:prelims}. This allows us to make connections with the control as inference perspective \cite{levine2018reinforcement} as it applies to the constrained reinforcement learning task formalized in Eq.~\ref{eq:RL}. Suppose we are given a model parameterized by $\pi_{\theta}(y\mid x)$, and the goal is to minimize the KL divergence $\mathbb{D}_{\text{KL}}[\pi_{\theta}(y|x) \mid \mid \pi^*(y\mid x)]$, where $\pi^*$ is the optimal policy from Eq.~\ref{eq:optimum_model} as determined by the reward function $r_{\phi}(y, x)$. Basic manipulation then yields the following objective function:

\begin{equation}\label{eq:AC}
    \max_{\pi_{\theta}}\mathbb{E}_{\pi_{\theta}(y\mid x)}\bigg[\underbrace{r_{\phi}(x, y) -\beta\log\sum_{y}\pi_\text{ref}(y\mid x)\exp\left(\frac{1}{\beta}r_{\phi}(x, y)\right)}_{f(r_{\phi}, \pi_\text{ref}, \beta)} - \underbrace{\beta\log\frac{\pi_{\theta}(y\mid x)}{\pi_\text{ref}(y\mid x)}}_{\text{KL}}\bigg]
\end{equation}

This objective is identical to that optimized in previous studies 
\citep{ziegler2020finetuning, stiennon2022learning, bai2022training, ouyang2022training}, where the DPO-equivalent reward is applied within the $r_{\phi}$ reward class. In this framework, the normalization factor within $f(r_{\phi}, \pi_\text{ref}, \beta)$ can be viewed as the soft value function associated with the reference policy $\pi_\text{ref}$. Although this normalization does not impact the final optimal policy, omitting it may result in high-variance policy gradients, potentially destabilizing the learning process. One way to account for the normalization is through a learned value function; however, this presents its own optimization challenges. Alternatively, prior works have often employed reward normalization by referencing a human completion as a baseline, which essentially provides a single-sample Monte-Carlo estimate of the normalizer. The DPO reparameterization, by comparison, defines a reward function that eliminates the need for such baselines.

\section{Experiments}
Here, we conduct empirical studies to investigate how well DPO learns policies from preference data. We begin with a controlled text-generation environment and examine DPO’s sample efficiency in balancing reward maximization and KL-regularization against the reference policy, in comparison with prevalent preference learning algorithms such as PPO. Subsequently, we apply DPO to larger models and more complex RLHF benchmarks, including summarization and dialogue tasks. Across these tasks, we observe that DPO, with minimal hyperparameter tuning, consistently matches or surpasses the performance of established baselines like PPO-based RLHF, as well as approaches that select the best of $N$ trajectories under a learned reward. The experimental setup is described below, and further implementation details can be found in Appendix~\ref{app:exp_details}.

\textbf{Tasks.} We investigate three distinct tasks involving open-ended text generation. In every case, the learning algorithms derive a policy using a preference dataset $\mathcal{D}=\bigl\{x^{(i)}, y_w^{(i)}, y_l^{(i)}\bigr\}_{i=1}^N$. For \textbf{controlled sentiment generation}, the input $x$ comprises the initial segment of a movie review selected from the IMDb dataset \cite{maas-EtAl:2011:ACL-HLT2011}, and the goal is to produce a continuation $y$ exhibiting positive sentiment. To facilitate systematic evaluation, we generate pairs of outputs by sampling generations and leveraging a pre-trained sentiment classifier to establish preferences, such that $p(\text{positive}\mid x,y_w)>p(\text{positive}\mid x,y_l)$. For the SFT baseline, GPT-2-large is fine-tuned on IMDB training reviews until performance plateaus (additional information provided in App~\ref{app:sentiment_details}). In the \textbf{summarization} task, the input $x$ is drawn from Reddit forum posts, with the policy required to produce a summary $y$ capturing the post’s main ideas. Consistent with earlier studies, we employ the Reddit TL;DR summarization corpus \citep{volske-etal-2017-tl} along with human preference data obtained by \citeauthor{stiennon2022learning}. The corresponding SFT model is trained on human-authored Reddit summaries\footnote{\url{https://huggingface.co/CarperAI/openai_summarize_tldr_sft}}, using the TRLX framework \citep{leandro_von_werra_2023_7790115} for RLHF experiments. It should be noted that human preference labels are sourced from samples produced by a different but comparably fine-tuned SFT model, as curated by \citeauthor{stiennon2022learning}. Lastly, the \textbf{single-turn dialogue} task considers $x$ as a user prompt, ranging from technical questions to personal advice requests. The policy’s objective is to craft a useful and engaging reply $y$. Here, we utilize the Anthropic Helpful and Harmless dialogue dataset \citep{bai2022training}, which comprises 170,000 human-assistant dialogues, with each session concluding with two model-generated responses and a human-indicated preferred completion. Since there is no pre-existing SFT model for this dataset, we fine-tune a publicly available language model solely on the preferred completions to establish an SFT baseline.

\textbf{Evaluation.} Our evaluation methodology encompasses two main strategies. In the controlled sentiment generation scenario, to assess how effectively each method maximizes constrained rewards, we characterize their performance by plotting the set of achievable trade-offs between reward and KL-divergence from a reference policy—referred to as the frontier. The feasibility of this analysis relies on access to the actual reward function, provided here by a sentiment classifier. Conversely, practical applications often lack an explicit reward signal. For such real-world cases, we compare methods using their \textit{win rate} relative to a baseline policy, leveraging GPT-4 as an automatic surrogate for human raters. Specifically, GPT-4 is tasked with judging either summary quality in the summarization task or response helpfulness in the single-turn dialogue task. In these settings, for summarization, reference summaries from the test corpus serve as the baseline, whereas, for dialogue, the ground-truth preferred responses play that role. While literature indicates that language models may surpass traditional metrics as evaluators \citep{Chen2023ExploringTU}, we corroborate our choice of GPT-4 through a human subject experiment, discussed in Sec.~\ref{sec:human-judgments}. Our findings show that GPT-4 evaluations are highly consistent with human judgments, with agreement rates that match or exceed those observed among human annotators.

\textbf{Methods.} Beyond DPO, our analysis encompasses multiple established techniques for aligning language model outputs with human preferences. For summarization, we assess a zero-shot prompting protocol with \textbf{GPT-J} \citep{gpt-j}; for dialogue, we utilize a 2-shot prompting approach with \textbf{Pythia-2.8B} \citep{biderman2023pythia}. Additionally, we benchmark the \textbf{SFT} model and a \textbf{Preferred-FT} model, where the latter is obtained by supervised fine-tuning on the preferred output $y_w$, drawing $y_w$ either from SFT predictions (for sentiment and summarization) or from a generic language model (for dialogue). We also include \textbf{Unlikelihood}~\citep{welleck2019neural}, a pseudo-supervised approach that promotes high probability for $y_w$ and simultaneously penalizes likelihood assigned to the dispreferred completion $y_l$, modulated by an optional coefficient $\alpha\in[0,1]$ on the 'unlikelihood' component. Furthermore, we consider \textbf{PPO} \citep{schulman2017proximal}, where the policy is trained on a reward function inferred from preference labels, alongside \textbf{PPO-GT}, which leverages the actual ground-truth reward during training in the sentiment-controlled environment. For sentiment tasks, we evaluate two PPO-GT variants: a standard implementation \cite{leandro_von_werra_2023_7790115} and a modified version where rewards are normalized and hyperparameters adjusted for optimal results (these modifications also apply to PPO models trained with learned rewards). Lastly, the \textbf{Best of $N$} baseline is included; this approach draws $N$ samples from either the SFT model (or Preferred-FT in dialogue) and selects the highest-reward response using a reward predictor trained on preference annotations. While this method often attains strong empirical performance and decouples reward model efficacy from PPO's optimization dynamics, it is resource-intensive, as it necessitates generating $N$ responses for every test prompt.

\subsection{How well can DPO optimize the RLHF objective?}

In standard RLHF methods, the reward maximization objective constrained by KL divergence ensures a trade-off between maximizing the reward and keeping the learned policy close to a reference policy. Thus, a meaningful evaluation of different approaches must jointly consider the attained reward and the corresponding KL divergence; obtaining marginally higher rewards at the cost of greatly increased KL is not an ideal outcome. Figure~\ref{fig:frontier-tldr-main} presents the trade-off frontier between reward and KL across several methods within the sentiment domain. For each algorithm, we conduct a series of training sessions, varying a hyperparameter that controls policy conservativeness (target KL values of $\{3,6,9,12\}$ for PPO, $\beta$ values of $\{0.05,0.1,1,5\}$, $\alpha$ values of $\{0.05,0.1,0.5,1\}$ for unlikelihood, and random seeds for preferred-FT), comprising a total of 22 experiments. At intervals of 100 training steps, extending through convergence, we evaluate policies using a dedicated test prompt set. For each policy, we calculate both the mean reward based on the ground truth reward signal and the average sequence-level KL divergence\footnote{Specifically, this refers to the aggregate sum of per-timestep KL-divergences.} from the reference policy $\text{KL}\left(\pi\mid \mid \pi_\text{ref}\right)$. Our findings indicate that DPO consistently yields the most effective frontier, achieving superior rewards with comparatively low KL values. This finding is significant for several reasons. Primarily, although DPO and PPO are designed to optimize the same reward-KL objective, DPO demonstrates a clear advantage in efficiency, with its reward-to-KL trade-off dominating that of PPO. Furthermore, DPO outperforms PPO's frontier \emph{even in scenarios where PPO is provided access to the exact reward function} (PPO-GT).

\subsection{Can DPO scale to real preference datasets?}
\label{sec:dpo-real-datasets}
We proceed to assess the fine-tuning capabilities of DPO on tasks involving summarization and single-turn dialogue. In the context of summarization, widely used automated metrics like ROUGE often show weak alignment with actual human judgments~\citep{stiennon2022learning}, and past studies have demonstrated that fine-tuning language models with PPO based on human feedback yields more compelling summaries. To benchmark the various approaches, we generate samples on the TL;DR summarization dataset’s test split, then calculate the mean win rate when these samples are compared to reference outputs in the test set. Completions from all evaluated methods are generated across a range of sampling temperatures from 0.0 to 1.0; the resulting win rates are depicted in Figure~\ref{fig:frontier-tldr-main} (right). DPO, PPO, and Preferred-FT methods all fine-tune using the identical GPT-J SFT model\footnote{\url{https://huggingface.co/CarperAI/openai_summarize_tldr_sft}}. Experimental results indicate that DPO attains a win rate of around 61\% at a temperature of 0.0, surpassing PPO, which achieves roughly 57\% at its optimal (also 0.0) temperature setting. Furthermore, DPO demonstrates a higher peak win rate than the best of $N$ baseline. It is important to point out that we did not perform significant tuning of the $\beta$ hyperparameter for DPO, so these findings may not reflect its full potential. In addition, we observe that DPO maintains performance more consistently across different sampling temperatures compared to PPO, whose effectiveness declines toward that of the original GPT-J model as temperature increases. The Preferred-FT approach, in contrast, shows minimal improvement over the SFT model baseline. Further, we present a direct comparison of DPO and PPO using human evaluators in Section~\ref{sec:human-judgments}, where DPO outputs sampled at a temperature of 0.25 were favored 58\% of the time relative to PPO outputs at the same temperature.

For single-turn dialogue, we assess the various approaches on a portion of the Anthropic HH dataset's test split \citep{bai2022training} that features a single exchange between a human and an assistant. To evaluate using GPT-4, the preferred responses from the test set serve as references when calculating the win rates across methods. Since no established SFT model exists for this particular task, we initialize with a pre-trained Pythia-2.8B model, employ Preferred-FT to fine-tune a reference model on selected completions—ensuring outputs align with the model's distribution—and subsequently apply DPO for training. Our comparison set also includes the best completion among 128 Preferred-FT generations (noting that the Best of $N$ method’s performance plateaus at 128, as detailed in Appendix Figure~\ref{fig:best-of-n}), as well as a 2-shot prompted Pythia-2.8B base model; under optimal temperature settings, DPO achieves performance that matches or surpasses these alternatives. In addition, we test an RLHF model trained via PPO on the Anthropic HH dataset\footnote{\url{https://huggingface.co/reciprocate/ppo_hh_pythia-6B}}, sourced from a reputable implementation\footnote{\url{https://github.com/CarperAI/trlx/tree/main/examples/hh}}, but do not identify any prompt or temperature configuration where this model outperforms the base Pythia-2.8B. Given our findings from TL;DR and the shared reward objective of both strategies, we treat Best of 128 as an approximate benchmark for PPO-based performance. Taken together, DPO stands out as the only method that is both computationally efficient and consistently exceeds the preferred completion baseline on the Anthropic HH dataset, offering performance comparable to or better than the more computationally intensive Best of 128 approach. Finally, as depicted in Figure~\ref{fig:dialogue-main}, DPO achieves its optimal results with relatively rapid convergence.

\subsection{Generalization to a new input distribution}

To more thoroughly assess the robustness of PPO and DPO to changes in input data distributions, we tested the policies learned from our Reddit TL;DR summarization task on an out-of-distribution dataset: news articles from the CNN/DailyMail test set \citep{nallapati-etal-2016-abstractive}. We applied the optimal sampling temperatures identified for TL;DR (0 and 0.25) and summarized the outcomes in Table~\ref{tab:ood}. To measure summary quality, we used GPT-4 win rates versus ground-truth summaries from the datasets, adopting the same GPT-4 (C) evaluation prompt as in the Reddit TL;DR setup, but substituting ``forum post'' for ``news article.'' On this new data distribution, the DPO policy still surpassed PPO by a wide margin. These results suggest that DPO models maintain generalization capabilities on par with PPO, despite not leveraging the extra unlabeled Reddit TL;DR prompts available to PPO.

\subsection{Validating GPT-4 judgments with human judgments}
\label{sec:human-judgments}
To assess the validity of GPT-4's summary evaluations, we conducted a human study focusing on the TL;DR summarization setting and employing two distinct prompts to GPT-4. The \textbf{GPT-4 (S)} (simple) prompt directs GPT-4 to identify which summary more effectively captures the essential information from the post. In contrast, the \textbf{GPT-4 (C)} (concise) prompt further requests a judgment on which summary is more succinct; we include this because, using \textbf{GPT-4 (S)}, GPT-4 tends to favor longer and more repetitive outputs than humans do. The full prompt texts are given in Appendix~\ref{app:prompts}. Three sets of comparisons were performed: between DPO (temp. 0.25), PPO (temp. 1.0), and SFT (temp. 0.25)—spanning a range of summary quality—against greedily-decoded PPO outputs, which correspond to the most successful PPO temperature setting. We observe that across both prompts, GPT-4 aligns with human raters roughly as often as humans concur with each other, indicating that GPT-4 can reliably stand in for human evaluation in this context (as the number of human annotators was limited, multiple human ratings were only available for DPO and PPO-1 comparisons). The \textbf{GPT-4 (C)} prompt produces win rates that more closely mirror human preferences and is thus adopted for the key results reported in Section~\ref{sec:dpo-real-datasets}. Additional details about the human evaluation protocol, user interface, and participant list are found in Appendix~\ref{app:human-study}.

\section{Discussion}
Preference-based learning provides a robust and adaptable approach for developing effective and aligned language models. In this work, we presented DPO, a streamlined training strategy that leverages preferences to guide language model learning, entirely sidestepping the need for reinforcement learning techniques. Unlike methods that fit the preference learning task within a reinforcement learning framework for the sake of leveraging conventional RL algorithms, DPO establishes a direct correspondence between language model policies and reward functions. This approach facilitates preference alignment using a straightforward cross-entropy loss, eliminating the reliance on reinforcement learning and preserving generality. Remarkably, DPO achieves performance on par with or superior to current RLHF approaches, such as those employing PPO, without the necessity for extensive hyperparameter adjustment, significantly lowering the entry barrier for preference-driven language model training.

\textbf{Limitations \& Future Work.} Our findings highlight several open questions meriting further investigation. For instance, how does the DPO-trained policy fare in generalizing beyond its training distribution, particularly relative to approaches that optimize against an explicit reward function? Preliminary results indicate DPO’s generalization capabilities may be comparable to models trained with PPO, yet broader evaluation is warranted. Further, it remains to be seen whether self-labeling mechanisms, when applied within the DPO framework, can effectively utilize unlabeled input prompts. Another pertinent line of inquiry concerns the manifestation of reward over-optimization in direct preference optimization—whether the modest performance dip observed in Figure~\ref{fig:dialogue-main}-right is indicative of this phenomenon. Though our evaluation extends to models with up to 6B parameters, expanding DPO to much larger, state-of-the-art models presents an exciting avenue for future study. In terms of evaluation, we observed that GPT-4’s computed win rates are sensitive to the specific prompt employed; refining automated assessment techniques to yield more consistent and high-quality judgments could be a productive area of research. Lastly, the potential applications for DPO span beyond human preference training in language models, offering possibilities in training generative models across a range of modalities.